\documentclass[11pt]{article}
\usepackage[margin=25mm]{geometry}
\usepackage{fontspec}
\setmainfont{Latin Modern Roman}
\usepackage{amsmath,amssymb,mathtools}
\usepackage{unicode-math}
\setmathfont{Latin Modern Math}
\usepackage{xurl,hyperref}
\hypersetup{colorlinks=true,urlcolor=blue}
\setcounter{secnumdepth}{0}
\setlength{\parindent}{0pt}
\setlength{\parskip}{5pt}
\setlength{\emergencystretch}{3em}
\newcommand{\sref}[2]{\hyperlink{source:#1:#2}{[S#2]}}
\newcommand{\cataloguescope}[1]{\paragraph{Recorded target.}#1}
\newcommand{\R}{\mathbb R}
\newcommand{\C}{\mathbb C}
\newcommand{\N}{\mathbb N}
\newcommand{\Z}{\mathbb Z}
\newcommand{\Q}{\mathbb Q}
\newcommand{\F}{\mathbb F}
\newcommand{\eps}{\varepsilon}
\begin{document}
% UnsolvedMath ID 30005988; source catalogue code OWR-14298586-004
\setcounter{section}{30005987}
\section{Cloud-Organization Limit Cycles}\label{30005988}
{\small UnsolvedMath ID 30005988 · Mathematical Physics}
\subsection{Definitions and mathematical statement}

\paragraph{Shared notation.}$\mathbb N=\{1,2,\ldots\}$; the standard number fields have their usual meanings. All additional conventions are stated locally.


A limit cycle is an isolated periodic orbit of a dynamical model of cloud organization. Its size, period and amplitude describe different observables. Quasistatic response means tracking the family of equilibrated or periodic states under slow forcing. Diurnal radiation and open boundaries alter the forcing and exchange of mass and energy.

\paragraph{Catalogue wording retained for traceability.}
Explain the mechanism producing limit cycles in shallow cumulus organization. Determine how external forcings affect their size, frequency, and amplitude; whether they respond quasi-statically to time-dependent forcing; how the diurnal radiative cycle and realistic open boundaries modify them; and how relevant they are to cloud feedback.

\paragraph{Scope and source qualification.}All forcing-response questions are retained, including whether time-dependent forcing is followed quasistatically.
\subsection{Short English statement}
Explain cloud-organization cycles and how forcing, daily radiation and open boundaries change their scales and climate feedback.
\subsection{Research attempt}
\paragraph{Model, date and investigated scope.}
OpenAI GPT-6 Astra (configured model \texttt{gpt-6-astra}), ultra reasoning; generated 11 October 2026. Solved a cycle normal form and computed its exact linear periodic forcing response, identifying a quantitative failure of quasistatic tracking near onset.

\paragraph{Sources and literature check.}
Read the full supplied definition, catalogue wording, source qualifications and listed references for 30005988. No fresh status or novelty search was performed. The argument below proves its auxiliary statements directly and does not import an unverified external theorem.

\paragraph{A cycle-response benchmark.} Use the auxiliary complex normal form \(\dot z=(\mu+i\omega)z-|z|^2z\), with \(\mu>0\), \(\omega\ne0\). Writing \(z=re^{i\theta}\) gives \(\dot r=\mu r-r^3\), \(\dot\theta=\omega\). The circle \(r_* =\sqrt\mu\) is attracting because the radial linearization is \(-2\mu\), and its period is \(2\pi/|\omega|\). Amplitude, frequency and radial recovery time are already distinct quantities in this minimal model.
\paragraph{Failure of automatic quasistatic tracking.} Let \(\mu(t)=\mu_0+\delta\cos(\Omega t)\), \(\mu_0>0\), with small \(\delta\), and write \(r=\sqrt{\mu_0}+\delta q+O(\delta^2)\). The linear response solves
\[
 \dot q+2\mu_0q=\sqrt{\mu_0}\cos(\Omega t).
\]
Its periodic solution has amplitude \(\sqrt{\mu_0}/\sqrt{4\mu_0^2+\Omega^2}\) and lag \(\arctan(\Omega/(2\mu_0))\). The quasistatic prediction instead gives amplitude \(1/(2\sqrt{\mu_0})\) and zero lag. Agreement therefore requires \(|\Omega|\ll2\mu_0\), and this requirement becomes restrictive near cycle onset.
\paragraph{Source comparison.} No reduction of the source cloud model to this normal form is established. In particular a diurnal forcing period need not be slow relative to the actual restoring time, and open boundaries may change the existence of the cycle itself. The first gap is extracting and validating dynamical coefficients and fluxes from the cloud model, including forcing-dependent cloud feedback. Mechanism, spatial size, amplitude, frequency, realistic boundaries and climate relevance all remain unresolved by this benchmark.

\paragraph{Verification and adversarial audit.}
Checked polar equations, attracting radial eigenvalue, response amplitude and phase, zero-frequency limit and loss of a uniform slow-forcing condition as mu tends to zero. Model self-audit only; no independent human mathematical verification. Applied contribution prompt A before drafting and B afterwards. Checked polar equations, attracting radial eigenvalue, response amplitude and phase, zero-frequency limit and loss of a uniform slow-forcing condition as mu tends to zero. First remaining gap: Derive and validate the cloud-model cycle mechanism and response coefficients, including open-boundary fluxes and cloud-feedback observables. The inherited source status is preserved; original-target completion remains zero.

Coordinator checked the radial cycle stability and forced linear response amplitude and phase lag. This additional check was model-assisted, not independent human verification.

\paragraph{Tools, process and contributions.}
Python 3 and local file tools used to read definitions and serialize manuscripts. Calculations are symbolic reasoning; no numerical experiment or theorem prover. Contribution prompt A was followed by prompt B. GPT-6 Astra with ultra reasoning drafted the research and performed its stated self-audit; any additional coordinator contribution is identified in the verification and audit. A GPT-6-based Codex coordinator assembled the manuscript and checked repository integration; its exact variant and reasoning configuration were not exposed. No human mathematical contribution, independent human proof check or formal proof certification is supplied. Complete derivations and the consulted source references are included here; private reasoning transcripts are not published.

\paragraph{Final status and first remaining gap.}
\textbf{UNRESOLVED}. Derive and validate the cloud-model cycle mechanism and response coefficients, including open-boundary fluxes and cloud-feedback observables.

\section{Completion Estimate}
\noindent\textbf{COMPLETION: 0\%}
This is an elementary auxiliary calculation or an explicit test of a restricted strategy. No new portion of the full source target is established.

\subsection{Sources}
\begin{itemize}
\item[\hypertarget{source:30005988:1}{S1}] ULAM.AI and the UnsolvedMath Contributors, UnsolvedMath, record 30005988 (OWR-14298586-004). Catalogue curation under CC BY 4.0; definitions and English exposition added locally.
\url{https://huggingface.co/datasets/ulamai/UnsolvedMath}.
\item[\hypertarget{source:30005988:2}{S2}] Model Hierarchies in Atmosphere, Ocean, and Climate Sciences (2024), Oberwolfach Reports, DOI 10.4171/owr/2024/31, pp. 1731--1786. Original report and mathematical context. Incorporated report material remains under the source CC BY-SA 4.0 terms unless otherwise noted.
\url{https://doi.org/10.4171/owr/2024/31}.
\end{itemize}
\label{end:30005988}
\end{document}
